
\section{\trEffectiveCartierDivisor{}「＝」\trInvertible\trIdealSheaf}

我们现在给出描述\trScheme{}上\trInvertibleSheaf{}的另一条途径。相较\trWeilDivisor，这样做的一个优点是：它能在处处非\trReduced{}的\trScheme{}上给出\trLineBundle（这样的概形在任何余维数 $1$ 的\trPrimeIdeal{}处都不可能\trRegular）。

设 $\mathrm{D}\hookrightarrow X$ 是一个\trClosedSubscheme，其相应的\trIdealSheaf{} $\mathscr{I}$ 是一个\trInvertibleSheaf。于是 $\mathscr{I}$ 局部平凡；设 $U$ 是一个供作\trTrivialization{}的仿射\trOpenSet{} $\Spec A$。于是，\trClosedSubscheme{}正合列 \eqref{eq:closed-subscheme-exact-sequence}
\[
0\longrightarrow\mathscr{I}\longrightarrow\OO_X\longrightarrow\OO_{\mathrm{D}}\longrightarrow 0
\]
对应于
\[
0\longrightarrow I\longrightarrow A\longrightarrow A/I\longrightarrow 0
\]
其中 $I\cong A$（作为 $A$-\trModule）。于是 $I$ 由单个元素（设为 $a$）生成，而该正合列开始为
\[
0\longrightarrow A\xrightarrow{\;\times a\;}A.
\]
由于乘以 $a$ 是单射，故 $a$ 不是\trZerodivisor。我们得出结论：$\mathrm{D}$ 由单个非\trZerodivisor{}的方程局部切出。这正是 \S9.5.1 中给出的\emph{\trEffectiveCartierDivisor}的定义。这一论证显然可以反转，于是我们便有了一个关于\trEffectiveCartierDivisor{}的快捷新定义（一个作为\trInvertibleSheaf{}的\trIdealSheaf{} $\mathscr{I}$——或者，等价地，其相应的\trClosedSubscheme）。

\textbf{15.6.A.（容易的练习）}\quad 证明：$a$ 在相差一个可逆函数倍乘的意义下唯一。

在 $X$ 局部诺特的情形，我们可以使用\trAssociatedPoint{}的语言（\S6.6.2），于是可将这一定义重述为：$\mathrm{D}$ 由单个方程局部切出，而该方程在 $X$ 的任何\trAssociatedPoint{}处都不消失。

现在我们来定义一个对应于 $\mathrm{D}$ 的\trInvertibleSheaf。看似显然的定义是取 $\mathscr{I}_{\mathrm{D}}$，但我们却做「相反」的事。

\textbf{15.6.1.}\ \emph{\trhDefinition{}：对应\trEffectiveCartierDivisor{}的\trInvertibleSheaf。}\quad 若 $\mathrm{D}$ 是一个\trEffectiveCartierDivisor，则定义\trInvertibleSheaf{} $\OO(\mathrm{D})$ 为 $\mathscr{I}_{\mathrm{D}}^{\vee}$。（取对偶的理由见\mbox{练习~15.6.C}。）定义 $\OO(n\mathrm{D})$ 为 $\OO(\mathrm{D})^{\otimes n}$。特别地，定义 $\OO(-\mathrm{D})$ 为 $\OO(\mathrm{D})^{\vee}$，亦即 $\OO(-\mathrm{D})=\mathscr{I}_{\mathrm{D}}$。

\textbf{15.6.B.（练习）}\quad 若 $X$ 正规，检验这与我们先前的 $\OO(\mathrm{D})$ 定义——\mbox{重要定义~15.4.6}——相一致。（在\mbox{定义~15.4.6} 中，$\mathrm{D}$ 是一个\trWeilDivisor；而在这道练习中，$\mathrm{D}$ 是一个\trEffectiveCartierDivisor。所以你得先在正规\trNoetherian{}概形 $X$ 上定义出对应于一个\trEffectiveCartierDivisor{}的\trWeilDivisor。）

由于\trIdealSheaf{} $\mathscr{I}_{\mathrm{D}}$ 就是 $\OO(-\mathrm{D})$，\trClosedSubscheme{}正合列 \eqref{eq:closed-subscheme-exact-sequence} 可以写作
\[
0\longrightarrow\OO(-\mathrm{D})\longrightarrow\OO\longrightarrow\OO_{\mathrm{D}}\longrightarrow 0.
\]

\trInvertibleSheaf{} $\OO(\mathrm{D})$ 带有一个典范\trSection{} $s_{\mathrm{D}}$：把 $0\to\mathscr{I}\to\OO$ 与 $\mathscr{I}^{\vee}$ 作张量积，便给出 $\OO\to\mathscr{I}^{\vee}$。（一个容易的不重要事实：我们本来也可以不对 $\mathscr{I}\to\OO$ 与 $\mathscr{I}^{\vee}$ 作张量积，而是对偶化 $\mathscr{I}\to\OO$，并且同样会得到这个\trSection。）

\textbf{15.6.C.（重要而出奇刁钻的练习）}\quad 回想：$X$ 上\trLocallyFreeSheaf{}的一个\trSection{}会切出 $X$ 的一个\trClosedSubscheme（\mbox{练习~14.2.D}）。证明：\trSection{} $s_{\mathrm{D}}$ 切出 $\mathrm{D}$。（与\mbox{注记~15.4.10} 相比较。）

这个构造是可逆的：

\textbf{15.6.D.（练习）}\quad 设 $\mathscr{L}$ 是一个\trInvertibleSheaf，而 $s$ 是一个局部地非\trZerodivisor{}的\trSection。（把这句话的含义弄清楚！特别地，若 $X$ 局部诺特，它的意思是「$s$ 不在 $X$ 的任何\trAssociatedPoint{}处消失」；见 \S6.6.2。）证明：$s=0$ 切出一个\trEffectiveCartierDivisor{} $\mathrm{D}$，且 $\OO(\mathrm{D})\cong\mathscr{L}$。

\textbf{15.6.E.（练习）}\quad 设 $\mathscr{I}$ 与 $\mathscr{J}$ 是两个可逆\trIdealSheaf（从而分别对应于\trEffectiveCartierDivisor{} $\mathrm{D}$ 与 $\mathrm{D}'$）。证明：$\mathscr{I}\mathscr{J}$ 是一个可逆\trIdealSheaf。（我们定义\textbf{两个\trQC\trIdealSheaf{}的积} $\mathscr{I}\mathscr{J}$，一如你所期望的：在每个仿射上，取两个相应\trIdeal{}的积——正如\mbox{练习~3.4.C} 中所定义的那样。为确保这一定义\trWellDefined，我们只需检验：若 $A$ 是一个环，$f\in A$，而 $I$、$J\subset A$ 是两个\trIdeal，则在 $A_f$ 中有 $(IJ)_f=I_fJ_f$。）我们把相应的\trCartierDivisor{}定义为 $\mathrm{D}+\mathrm{D}'$。验证：$\OO(\mathrm{D}+\mathrm{D}')\cong\OO(\mathrm{D})\otimes\OO(\mathrm{D}')$。

这样，我们便在\emph{\trEffectiveCartierDivisor}（其\trIdealSheaf{}可逆的\trClosedSubscheme；等价地说，由单个\trNonZerodivisor{}局部切出；或者在局部诺特的情形，由单个在任何\trAssociatedPoint{}处都不消失的方程局部切出）与\emph{有序对}（ordered pairs）$(\mathscr{L},s)$——其中 $\mathscr{L}$ 是一个\trInvertibleSheaf，而 $s$ 是一个局部地非\trZerodivisor{}的\trSection（或者在局部诺特的情形，不在任何\trAssociatedPoint{}处消失）——之间有了一个重要的对应。\trEffectiveCartierDivisor{}构成一个\trAbelianSemigroup。我们有一个半群之间的映射：从\trEffectiveCartierDivisor{}走到带有一个局部非\trZerodivisor{}之\trSection{}的\trInvertibleSheaf（从而也走到\trInvertibleSheaf{}的\trPicardGroup）。

\textbf{15.6.2.}\ \emph{到\trEffectiveCartierDivisor{}的正规（线）丛。}\quad 若 $\mathrm{D}$ 是\trScheme{} $X$ 上的一个\trEffectiveCartierDivisor，则把\trInvertibleSheaf{} $\OO(\mathrm{D})$ 限制到 $\mathrm{D}$ 自身上所得的\trRestriction，常称为到 $\mathrm{D}$ 的\textbf{正规（线）丛}（normal (line) bundle）（尽管或许它应当被称作正规可逆层）。我们把它记作 $\mathscr{N}_{\mathrm{D}/X}:=\OO_X(\mathrm{D})\vert_{\mathrm{D}}$。采用这套语言的动机在于：若 $X$ 与 $\mathrm{D}$ 都是\trSmooth{}的\trComplex\trVariety，则这种代数几何意义下的\trNormalBundle{}与微分几何意义下的\trNormalBundle{}相一致。在更「奇异」的情势下这一解释会失效，但这种直观即便在那时依然有用。我们将在更一般的设定中定义\trNormalBundle{}（与正规层，normal sheaf），见 \S21.2.15。

\textbf{15.6.3.}\ \emph{不重要的注记。}\quad 取两个\trEffectiveCartierDivisor{}之差，我们便得到许多\trInvertibleSheaf。事实上，我们「通常能把它们尽数得到」——要描述有限型 $k$-概形上一个无法这样描述的\trInvertibleSheaf，是非常困难的。例如：若该概形\trRegular、甚至是\trFactorial{}的，则一个也没有（就\trFactoriality{}而言，基本上由\mbox{命题~15.4.16} 可得；而正则概形都是唯一分解的，这由\mbox{Auslander--Buchsbaum 定理~13.8.5} 可知）。\mbox{练习~16.2.F} 将蕴涵：若该概形射影，则一个也没有。在其余一切合理的情势下它都成立；见 [Gr-EGA, IV$_4$.21.3.4]。然而它并不总成立；最早的也是最好的一个例子出自\trKleiman{}之手；见 [Kl3]。

\textbf{15.6.4.}\ \extrahard\ \textbf{\trCartierDivisor。}\quad 有一个与之相关的概念：\emph{\trCartierDivisor}（Cartier divisor）。许多人认为这一概念必不可少，所以我想试着解释一下，为什么除了在此简短一谈之外，我们不去讨论它。\trCartierDivisor{}的确重要，但可以说，并未重要到我们此刻所进行的讨论的层级上。使用\trCartierDivisor{}涉及技术上的困难，换来的却是少得出奇的视角与结果上的简化\footnote{原书此处作（suprisingly），当为（surprisingly）之误，径改。——译者注}。（欢迎专家在页边以\trCartierDivisor{}的语言发展一段平行的讨论，来检验这一点。）为完整起见，这里给出一个定义。在一个\trScheme{} $X$ 上，我们定义一个 $\OO_X$-\trModule{} $\mathscr{M}_X$：它在仿射\trOpenSet{} $\Spec A$ 上的诸\trSection{}就是 $A$ 的\trTotalFractionRing{}的（诸）元素（\S6.6.36），并配有「显而易见」的\trRestrictionMap。若 $X$ 是\trIntegral{}的，则 $\mathscr{M}_X$ 是对应于\trFunctionField{} $K(X)$ 的\trConstantSheaf。我们有一个诱导映射 $\OO_X\hookrightarrow\mathscr{M}_X$。令 $\mathscr{M}_X^{*}$ 为（在\trSection{}的层面上）对应于可逆元的\trGroup\trSheaf，它诱导 $\OO_X^{*}\hookrightarrow\mathscr{M}_X^{*}$。于是，$X$ 的\textbf{\trCartierDivisor}定义为 $\Gamma(X,\mathscr{M}_X^{*}/\OO_X^{*})$（连同其作为一个\trAbelian\trGroup{}的结构）。\trCartierDivisor{}的诸多表面上的优点，实际上来自其他的结构。\trCartierDivisor{}在\trMorphism{}下易于\trPullback，但\trLineBundle（或者，等价地，\trInvertibleSheaf）亦是如此。\trEffectiveCartierDivisor{}同样十分令人愉悦且简单。\trCartierDivisor{}易于产生\trInvertibleSheaf，但「每个\trInvertibleSheaf{}都来自某个\trCartierDivisor」却并不\emph{总是}真的，[Kl3]。

